\documentclass{beamer}

\usepackage{beamerthemeSingapore}
\usepackage[shorthands=off,brazil]{babel}
\usepackage[T1]{fontenc}
% sem guillemets: "> fora do código e << >> na fonte monoespaçada
\usepackage[protrusion=false,expansion=false]{microtype}
\DisableLigatures[<,>]{encoding = T1, family = tt*}
\usepackage[utf8]{inputenc}
\usepackage{graphicx}
\usepackage{url}
\usepackage{xcolor}
\usepackage{booktabs}
\usepackage{ragged2e}
\usepackage{etoolbox}
\usepackage{minted}
\usepackage{array}

\newcommand{\tcb}[1]{\textcolor{blue}{#1}}
\newcommand{\tcr}[1]{\textcolor{red}{#1}}
\newcommand{\tcg}[1]{\textcolor{green!50!black}{#1}}

\setminted{
  fontsize=\footnotesize,
  breaklines=true,
  breakanywhere=true,
  frame=lines,
  framesep=2mm,
  baselinestretch=1.05,
  tabsize=2
}

\apptocmd{\frame}{}{\justifying}{}

\setbeamertemplate{navigation symbols}{}

\title{14 - PHP: cookies, sessões e login}
\author{Prof. Alex Kutzke}
\date{Novembro de 2026}

\begin{document}

% ==========================================
\begin{frame}
	\begin{center}
		\large{Setor de Educação Profissional e Tecnológica}\\
		Tecnologia em Análise e Desenvolvimento de Sistemas\\
		\vspace{1cm}
		\small{DS122 - Desenvolvimento de Aplicações Web 1}
	\end{center}
	\titlepage
\end{frame}

\begin{frame}{Objetivos da aula}
	Ao final desta aula você deve ser capaz de:
	\begin{enumerate}
		\item Explicar por que o HTTP precisa de cookies e sessões, e onde cada um
		      guarda os dados;
		\item Criar, ler e remover um cookie em PHP;
		\item Guardar e ler dados em \texttt{\$\_SESSION}, e encerrar uma sessão;
		\item Gravar senhas com \texttt{password\_hash} e conferi-las com
		      \texttt{password\_verify};
		\item Construir um login e proteger as páginas e os scripts de uma área
		      administrativa.
	\end{enumerate}

	\vspace{0.2cm}
	Código de exemplo:\\
	{\small\mbox{\texttt{https://gitlab.com/ds122-alexkutzke/ds122-login-app}}}
\end{frame}

% ==========================================
\section{Estado}
% ==========================================

\begin{frame}[fragile]{Demonstração: contador de visitas}
	\begin{minted}{html+php}
<?php
session_start();
$_SESSION["visitas"] = ($_SESSION["visitas"] ?? 0) + 1;
?>
<p>Visitas nesta sessão: <?= $_SESSION["visitas"] ?></p>
	\end{minted}

	\begin{itemize}
		\item Rede (\texttt{F12}): \texttt{Set-Cookie: PHPSESSID=...} na primeira
		      resposta, \texttt{Cookie: PHPSESSID=...} nas requisições seguintes;
		\item Numa janela anônima, o contador volta a 1;
		\item O número fica num arquivo do servidor, \texttt{sess\_} seguido do
		      identificador:
	\end{itemize}

	\begin{minted}{text}
visitas|i:2;
	\end{minted}
\end{frame}

\begin{frame}{Cookies e sessões}
	O HTTP é \textbf{sem estado}: para reconhecer o usuário, o servidor precisa
	que o cliente reenvie alguma informação em toda requisição.

	\vspace{0.3cm}
	\small
	\begin{tabular}{p{0.27\textwidth} p{0.28\textwidth} p{0.31\textwidth}}
		\toprule
		 & \textbf{Cookie} & \textbf{Sessão} \\ \midrule
		Onde ficam os dados & No navegador & No servidor \\
		O que trafega & Os próprios dados & Só o identificador, num cookie \\
		O usuário altera os dados & Sim & Não \\
		Uso típico & Preferências & Usuário autenticado \\
		Em PHP & \texttt{setcookie()}, \texttt{\$\_COOKIE} &
		\texttt{session\_start()}, \texttt{\$\_SESSION} \\
		\bottomrule
	\end{tabular}
\end{frame}

% ==========================================
\section{Cookies}
% ==========================================

\begin{frame}[fragile]{Criar, ler e remover um cookie}
	Criar, com nome, valor, expiração e caminho:
	\begin{minted}[startinline]{php}
setcookie("tema", "escuro", time() + 60 * 60 * 24 * 30, "/");
	\end{minted}

	Ler:
	\begin{minted}[startinline]{php}
$tema = $_COOKIE["tema"] ?? "claro";
	\end{minted}

	Remover, com uma data no passado e o mesmo caminho:
	\begin{minted}[startinline]{php}
setcookie("tema", "", time() - 3600, "/");
	\end{minted}

	\begin{itemize}
		\item \texttt{setcookie()} vem antes de qualquer saída HTML;
		\item o cookie só aparece em \texttt{\$\_COOKIE} na \textbf{próxima}
		      requisição.
	\end{itemize}
\end{frame}

\begin{frame}{O usuário pode alterar o cookie}
	No \texttt{F12}, troque o valor de \texttt{tema} por
	\texttt{<b>azul</b>} e recarregue a página.

	\vspace{0.3cm}
	\begin{itemize}
		\item Um cookie é um dado enviado pelo usuário, como um campo de
		      formulário: exiba-o com \texttt{htmlspecialchars()};
		\item Um cookie \texttt{admin=1} não protege nada, porque qualquer
		      visitante pode criá-lo;
		\item A informação de quem está autenticado fica na sessão.
	\end{itemize}
\end{frame}

% ==========================================
\section{Sessões}
% ==========================================

\begin{frame}[fragile]{\texttt{session\_start()} e \texttt{\$\_SESSION}}
	\begin{minted}{php}
<?php
session_start();

// Gravar
$_SESSION["usuario_nome"] = "Maria Silva";

// Ler, em qualquer página que chame session_start()
echo htmlspecialchars($_SESSION["usuario_nome"] ?? "visitante");
	\end{minted}

	\begin{itemize}
		\item Se a requisição traz o cookie \texttt{PHPSESSID}, o PHP carrega os
		      dados daquela sessão. Se não traz, cria uma sessão nova;
		\item Toda página que usa \texttt{\$\_SESSION} chama
		      \texttt{session\_start()} uma vez, antes de qualquer saída HTML.
	\end{itemize}
\end{frame}

\begin{frame}[fragile]{Opções do cookie da sessão: \texttt{admin/sessao.php}}
	\begin{minted}{php}
<?php
session_start([
    "cookie_httponly" => true,
    "cookie_samesite" => "Lax",
]);
	\end{minted}

	Com uma cópia do identificador, outra pessoa acessa o site como se fosse o
	usuário.

	\begin{itemize}
		\item \texttt{cookie\_httponly}: o JavaScript da página não lê o cookie;
		\item \texttt{cookie\_samesite} com \texttt{"Lax"}: o navegador não envia
		      o cookie num formulário POST que outro site submeta para o seu.
	\end{itemize}
\end{frame}

% ==========================================
\section{Senhas}
% ==========================================

\begin{frame}[fragile]{\texttt{password\_hash} e \texttt{password\_verify}}
	O banco guarda o \emph{hash} da senha, do qual não se recupera o texto
	original.

	\begin{minted}[startinline]{php}
echo password_hash("123456", PASSWORD_DEFAULT), "<br>";
echo password_hash("123456", PASSWORD_DEFAULT), "<br>";
	\end{minted}

	\begin{minted}[fontsize=\scriptsize]{text}
$2y$10$q6WA7ddvt4RPkToscD5YPO6OjWIV1oNZaxPvxCrstTStCIi9DXavq
$2y$10$zITCq58z8tNBapehdpRRwuHmEYOwBC5ivKRBfLw/vYF71oULWke9W
	\end{minted}

	\begin{itemize}
		\item A mesma senha produz hashes diferentes, por causa do \emph{salt}. A
		      comparação com \texttt{===} não funciona;
		\item \texttt{password\_verify(\$senha, \$hash)} devolve \texttt{true}
		      quando a senha confere;
		\item Não use \texttt{md5()} nem \texttt{sha1()} para senhas.
	\end{itemize}
\end{frame}

% ==========================================
\section{Login}
% ==========================================

\begin{frame}[fragile]{Estrutura do exemplo}
	\begin{minted}[breaklines=false]{text}
ds122-login-app/
|-- config.php               dados de conexão
|-- conexao.php              conexão com o MySQL
|-- criar_bd.php             cria a base de dados
|-- criar_tabela.php         cria a tabela de usuários
|-- cadastrar.php            cadastra o primeiro usuário
|-- index.php                página pública
`-- admin/
    |-- sessao.php           inicia a sessão
    |-- verifica_login.php   barra quem não está autenticado
    |-- login.php            formulário e processamento
    |-- dashboard.php        página protegida
    `-- logout.php           encerra a sessão
	\end{minted}
\end{frame}

\begin{frame}[fragile]{\texttt{conexao.php}}
	\begin{minted}[fontsize=\scriptsize]{php}
<?php
require_once __DIR__ . "/config.php";

try {
    $conn = mysqli_connect(DB_HOST, DB_USER, DB_PASS, DB_NAME);
    mysqli_set_charset($conn, "utf8mb4");
} catch (mysqli_sql_exception $e) {
    // Em produção, mostre uma mensagem genérica, sem detalhes do banco
    die("Erro de conexão: " . $e->getMessage());
}
	\end{minted}

	\begin{itemize}
		\item Desde o PHP 8.1, um erro do \texttt{mysqli} lança uma exceção, e o
		      script para com \texttt{Fatal error} se ninguém a tratar;
		\item Sem \texttt{?>} no final: linhas em branco depois dele seriam
		      enviadas antes dos cabeçalhos, e o login falharia com
		      \texttt{headers already sent}.
	\end{itemize}
\end{frame}

\begin{frame}[fragile]{\texttt{cadastrar.php}}
	\begin{minted}[fontsize=\scriptsize,startinline]{php}
// O banco guarda o hash, nunca a senha
$hash = password_hash($senha, PASSWORD_DEFAULT);

try {
    $stmt = mysqli_prepare($conn, "INSERT INTO usuarios (nome, email, senha) VALUES (?, ?, ?)");
    mysqli_stmt_bind_param($stmt, "sss", $nome, $email, $hash);
    mysqli_stmt_execute($stmt);
    echo "Usuário cadastrado com sucesso!";
} catch (mysqli_sql_exception $e) {
    // 1062 é o código do MySQL para valor repetido em coluna UNIQUE
    if ($e->getCode() === 1062) {
        echo "Já existe um usuário com este e-mail.";
    } else {
        echo "Erro ao cadastrar: " . $e->getMessage();
    }
}
	\end{minted}

	Não é um formulário público: qualquer visitante poderia criar um usuário e
	entrar na área administrativa.
\end{frame}

\begin{frame}[fragile]{\texttt{admin/login.php}: busca do usuário}
	\begin{minted}[fontsize=\scriptsize,startinline]{php}
$email = trim($_POST["email"] ?? "");
$senha = $_POST["senha"] ?? "";

// 1. Busca o usuário pelo e-mail
$stmt = mysqli_prepare($conn, "SELECT id, nome, senha FROM usuarios WHERE email = ?");
mysqli_stmt_bind_param($stmt, "s", $email);
mysqli_stmt_execute($stmt);
$resultado = mysqli_stmt_get_result($stmt);
$usuario = mysqli_fetch_assoc($resultado); // null quando o e-mail não existe
	\end{minted}

	O \texttt{conexao.php} está na pasta de cima:
	\begin{minted}[fontsize=\scriptsize,startinline]{php}
require_once __DIR__ . "/../conexao.php";
	\end{minted}
\end{frame}

\begin{frame}[fragile]{\texttt{admin/login.php}: conferência da senha}
	\begin{minted}[fontsize=\scriptsize,startinline]{php}
// 2. Confere a senha digitada com o hash guardado no banco
if ($usuario && password_verify($senha, $usuario["senha"])) {
    // 3. Troca o identificador e grava o usuário na sessão
    session_regenerate_id(true);
    $_SESSION["usuario_id"]   = $usuario["id"];
    $_SESSION["usuario_nome"] = $usuario["nome"];

    header("Location: dashboard.php");
    exit;
}

$erro = "E-mail ou senha inválidos.";
	\end{minted}

	A mensagem é a mesma para e-mail inexistente e senha errada, e não revela
	quais e-mails estão cadastrados.
\end{frame}

\begin{frame}[fragile]{Páginas protegidas}
	\texttt{admin/verifica\_login.php}:
	\begin{minted}{php}
<?php
require_once __DIR__ . "/sessao.php";

if (!isset($_SESSION["usuario_id"])) {
    header("Location: login.php");
    exit;
}
	\end{minted}

	Primeira linha de \texttt{admin/dashboard.php} e de toda página protegida:
	\begin{minted}{php}
<?php
require_once __DIR__ . "/verifica_login.php";
?>
	\end{minted}
\end{frame}

\begin{frame}[fragile]{O \texttt{exit} depois do redirecionamento}
	\texttt{header("Location: ...")} apenas acrescenta um cabeçalho. Sem o
	\texttt{exit}, o script continua e envia a página protegida inteira.

	\vspace{0.2cm}
	Demonstração: sem o \texttt{exit} no \texttt{verifica\_login.php},

	\begin{minted}{bash}
curl -i http://localhost:8000/admin/dashboard.php
	\end{minted}

	\begin{minted}[fontsize=\scriptsize]{text}
HTTP/1.1 302 Found
Location: login.php
...
    <h1>Painel</h1>
    <p>Esta página só aparece para usuários autenticados.</p>
	\end{minted}
\end{frame}

\begin{frame}[fragile]{Scripts que alteram dados}
	Esconder o link ou o formulário não protege o script que recebe os dados.
	Todo arquivo da pasta \texttt{admin/} começa com a verificação:

	\begin{minted}{php}
<?php
// admin/excluir.php, admin/salvar.php e os demais
require_once __DIR__ . "/verifica_login.php";
require_once __DIR__ . "/../conexao.php";

// só então lê $_POST ou $_GET e altera o banco
	\end{minted}

	As exceções são \texttt{sessao.php}, \texttt{login.php} e
	\texttt{logout.php}.
\end{frame}

\begin{frame}[fragile]{\texttt{admin/logout.php}}
	\begin{minted}{php}
<?php
require_once __DIR__ . "/sessao.php";

// 1. Apaga os dados da sessão
$_SESSION = [];

// 2. Remove o cookie do navegador
setcookie(session_name(), "", time() - 3600, "/");

// 3. Destrói a sessão no servidor
session_destroy();

header("Location: login.php");
exit;
	\end{minted}
\end{frame}

\begin{frame}[fragile]{Fixação de sessão (\emph{session fixation})}
	Com o \texttt{contador.php} e um identificador inventado:
	\begin{minted}[fontsize=\scriptsize]{bash}
curl -b "PHPSESSID=abc123" http://localhost:8000/contador.php
curl -b "PHPSESSID=abc123" http://localhost:8000/contador.php
	\end{minted}

	\begin{itemize}
		\item A saída mostra 1 e, depois, 2: na configuração padrão, o PHP aceita
		      um identificador escolhido pelo cliente;
		\item Se a vítima usar um identificador que o atacante conhece, e ele
		      continuar valendo após o login, o atacante entra na conta;
		\item \texttt{session\_regenerate\_id(true)} cria um identificador novo no
		      login e apaga o antigo.
	\end{itemize}
\end{frame}

% ==========================================
\section{Entrega 3}
% ==========================================

\begin{frame}{Do exemplo para a Entrega 3}
	O restante da aula é dedicado ao trabalho.

	\begin{enumerate}
		\item Crie a tabela \texttt{usuarios} e inclua o comando no arquivo
		      \texttt{.sql};
		\item Adapte os arquivos da pasta \texttt{admin/} à sua conexão e aos seus
		      nomes;
		\item Inclua \texttt{verifica\_login.php} na listagem, no cadastro, na
		      edição, na exclusão e na página de mensagens;
		\item Numa janela anônima, digite o endereço de cada arquivo de
		      \texttt{admin/} e confira que todos levam ao login;
		\item Teste a senha errada, o e-mail inexistente e o logout.
	\end{enumerate}

	\vspace{0.2cm}
	Prazo: 29/11 (domingo), às 23h59. Não há entrega em atraso.

	Próxima aula, em 04/12: Prova 3 e encerramento.
\end{frame}

\end{document}
